\section{Introduction}
Feldman, Kamath and Srebro ask whether distribution-independent statistical-query learning forces a common low-dimensional linear representation \citep{FKS26}. We answer the SQ question negatively, including its polynomial and probabilistic-dimension variants. Distribution independence requires one learner to work under every marginal; it does not prohibit that learner from adapting to the marginal through its queries.

A fixed linear representation must realize every example--hypothesis sign simultaneously. A learner can first ask where the unknown marginal places mass, then label a region on which surviving hypotheses agree or use a disagreement to eliminate hypotheses. Large monochromatic rectangles under every product distribution are precisely what a marginal-free mixture can cover pointwise, by finite minimax. We convert such a mixture into an adversarial-oracle learner and charge either action to the product of the rectangle's two masses. Random incidence flips have this rectangle structure but large sign-rank. Therefore their advantage over common linear representations is not inherently tied to learning under one fixed marginal.

\Needspace{5\baselineskip}
\paragraph{Contributions.}
\begin{itemize}
\item RECTIFY learns whenever $\rect(H)\ge1/\rho$, using $O(\rho(\log|H|+\log(1/\epsilon)))$ queries at tolerance $\Omega(\epsilon/(\rho\log(1/\epsilon)))$, with a worst-case query budget and expected-error guarantee against every valid adaptive oracle.
\item On the $2M$-point incidence-flip construction, fixed-accuracy learning uses $O(\log M)$ queries although $\dc=\Omega(M^{1/3}/\log M)$, $\dc_{1/2}=\Omega(M^{1/3}/\log^2M)$, and $\dc^\eta=\Omega_\eta(M^{1/3}/\log M)$ for fixed $\eta<1/2$. This refutes all three proposed polynomial implications.
\item Fixed-dimensional cap mixtures yield table-polynomial and deterministic implementations. Proper RECTIFY learns the same classes with $O(\log M+1/\epsilon+\log(1/\epsilon))$ queries.
\item Under subexponential PRF security, almost every polynomial-length key gives a succinctly evaluable class with $\dc\ge n^c/(30\log_2n)$ for any prescribed fixed $c\ge1$. This conditionally refutes the linear form and every prescribed polynomial degree with succinctly evaluable entries.
\item An order-median learner resolves the all-light geometric case in $O(n^2)$ queries. Given a key, it is polynomial-time and improper for every flip; a regular-star condition, satisfied by almost every PRF key, makes it proper with the same fixed-accuracy order. Hence restricting to polynomial-time learners does not rescue the linear query/tolerance form.
\item Bounded average deterministic communication under every product distribution implies SQ learning, with exponential loss in the communication cost. Cube halfspaces show that both this condition and large rectangle ratio are sufficient but unnecessary. The global projective-plane ratio is $\Theta(1/q)$, and joint spectral certificates give the accompanying query/tolerance lower bound.
\item We separate the original class-dependent benign-SGD question from uniform hidden-key learning. Unseen-label and PRF reductions obstruct the latter, while neither changes the architecture quantifier of the former. Source-program length alone also fails as a resource repair when computation is free.
\end{itemize}

\paragraph{Imported results and the new step.}
The real-grid flip construction is due to \citet{HHPTZ22}, and its counting uses Warren's theorem \citep{Warren68,AMY16}. The homogeneous-rectangle theorem \citep{APP05} only sharpens the constant to $2^{-15}$; an independent cap argument suffices for the separation. The rectangle-mixture minimax device already appears in the average-margin proof of \citet[Theorem~3.1]{HHPTZ22}. Our contributions are the learner and its product-progress analysis, proper and succinct implementations, probabilistic-dimension reductions, and the stated consequences. Feldman's SQ characterization is used only in Section~\ref{sec:rsd}. The new conditional argument uses the classical singly exponential existential-real decision theorem and the explicitly stated PRF assumption. Section~\ref{sec:related} audits the specific 2017 and 2026 comparisons.

\paragraph{Scope and verification.}
Query complexity does not charge class description or computation; table-polynomial and seed-polynomial implementations are distinguished throughout. The high-rank unconditioned selection is non-explicit, whereas the conditional selection has a negligible exceptional key set rather than an efficiently identified good-key sequence. The PRF conclusions depend on the stated security assumption. The proofs use the identified classical and recent imported theorems. CPU experiments have exact response and finite-coverage audits; they do not certify adversarial-search optimality or cryptographic security. The pinned Lean sources remain uncompiled. The original Gaussian-initialized, bias-free online-SGD implication remains open.
